\chapter{Boundary Diagnostics and Interpolation Baselines}
\label{app:diagnostic_ablation_forward_cfd}

\section{Boundary Diagnostics}
\label{sec:diagnostics}

The diagnostics below define what the final recipe does \emph{not} claim. They are single-seed or local-sweep checks under the final LES-pivot, $1024$-collocation, $20\,000$-iteration protocol unless otherwise noted, and are therefore used to set boundaries rather than to support Chapter~\ref{ch:results} main evidence.

\begin{table}[!htbp]
    \centering
    \caption{Boundary diagnostics for alternative recipe choices at $Re=10\,000$, $K=100$. Metrics are percentages unless noted. The multi-seed main baseline remains $5.71 \pm 0.12\%$ KE error.}
    \label{tab:boundary_diagnostics}
    \footnotesize
    \setlength{\tabcolsep}{3pt}
    \renewcommand{\arraystretch}{1.08}
    \begin{tabularx}{\textwidth}{p{2.6cm}p{4.6cm}p{3.2cm}X}
    \toprule
    \textbf{Diagnostic} & \textbf{Probe} & \textbf{Observed range} & \textbf{Reading} \\
    \midrule
    Multi-constraint AL
    & Add AL terms to NS momentum and/or continuity, with alternative GradNorm task sets.
    & KE $5.47$--$6.31$; div ratio $0.39$--$0.57$.
    & NS-momentum AL is not categorically unstable, but the effect is a divergence--accuracy trade-off rather than a validated improvement. \\
    \midrule
    Continuity-AL strength
    & Vary $\rho$ around the main $\rho=0.1$ setting: $0.03$, $0.30$, $1.00$.
    & KE $5.88$--$6.12$; div ratio decreases $0.45 \to 0.28$.
    & Larger $\rho$ tightens divergence but costs field accuracy; $\rho$ is a physics--data knob, not a source of improved reconstruction. \\
    \midrule
    Forcing recovery
    & Learn $A$, $k_f$, or both from near-zero initialisation ($A_{\rm init}=0.001$, $k_{f,\rm init}=0.01$).
    & Learned values collapse near $A=0.001$, $k_f=0.01$ while KE stays $5.74$--$5.85$.
    & Accurate velocity reconstruction does not certify recovery of the forcing parameters; structural identifiability is not established. \\
    \midrule
    Filtering vs.\ smoothing
    & Compare deployable filtering against offline bidirectional smoothing.
    & Filtering KE $5.71 \pm 0.12$; smoothing KE $5.74$.
    & Future-sensor access gives no clear gain in this single-seed diagnostic; filtering remains the deployable default. \\
    \bottomrule
    \end{tabularx}
\end{table}

\paragraph{Recipe implication.}
The final pipeline therefore keeps continuity-only AL with $\rho=0.1$, fixed forcing parameters, and filtering-mode CfC. These choices are conservative: they are supported by the $n=5$ main baseline, avoid non-deployable future-sensor access, and do not promote single-seed alternatives into recipe claims.

\paragraph{Detailed diagnostic anchors.}
\label{sec:multi_al_negative}
\label{sec:al_strength_appendix}
\label{sec:forcing_identifiability}
\label{sec:filtering_smoothing}
Table~\ref{tab:boundary_diagnostics} is the canonical summary for the diagnostic boundary cases referenced from Chapter~\ref{ch:results} and Chapter~\ref{ch:discussion_conclusion}.

\section{Engineering-Transferable Interpolation Baselines}
\label{app:fair_baselines}

Three classical reconstructions respect the same sensor-only constraint as PI-CON (no DNS access at training or inference): radial basis function (RBF) interpolation with a multiquadric kernel, inverse distance weighting (IDW), and a divergence-free trigonometric least-squares fit on a truncated Fourier basis ($k_{\max}=5$, matching the $K=100$ Nyquist scale). Table~\ref{tab:fair_baselines} compares them with PI-CON on the same $K=100$ LES-derived sensors used for the main baseline, so the comparison isolates the reconstruction method at a fixed sensor placement. The RBF shape parameter ($\varepsilon=10$, of order one sensor spacing) and the truncation ($k_{\max}=5$, the $K=100$ Nyquist scale) are fixed from these a-priori scales rather than tuned to minimise DNS error, keeping the baselines sensor-only.

\begin{table}[!htbp]
    \centering
    \caption{Engineering-transferable interpolation baselines vs.\ PI-CON at $Re=10\,000$, $K=100$ (same $K=100$ LES-derived sensors as the main baseline, no DNS access). PI-CON: mean $\pm$ standard deviation over $n=5$ seeds. Rightmost column: PI-CON relative reduction in $u$ rel-$L_2$. Relative-error metrics are percentages.}
    \label{tab:fair_baselines}
    \footnotesize
    \setlength{\tabcolsep}{3pt}
    \begin{tabular}{lccccc}
    \toprule
    \textbf{Method} & \textbf{KE \%} & \textbf{$u$ $L_2$ \%} & \textbf{$v$ $L_2$ \%} & \textbf{$\omega$ $L_2$ \%} & \textbf{$u$-$L_2$ rel.\ red.} \\
    \midrule
    RBF multiquadric ($\varepsilon=10$)   & $5.08$  & $30.02$ & $34.43$ & $58.33$ & $54.5\%$ \\
    IDW ($p=2$)                           & $66.66$ & $52.88$ & $62.02$ & $81.89$ & $74.2\%$ \\
    Div-free trig.\ LSQ ($k_{\max}=5$)    & $4.42$  & $25.87$ & $31.96$ & $63.41$ & $47.2\%$ \\
    \midrule
    \textbf{PI-CON (ours, $n=5$)}         & $5.71 {\pm} 0.12$ & $\mathbf{13.65 {\pm} 0.06}$ & $\mathbf{17.52 {\pm} 0.10}$ & $\mathbf{41.77 {\pm} 0.12}$ & -- \\
    \bottomrule
    \end{tabular}
\end{table}

\paragraph{Kinetic energy is a misleading sole metric.}
RBF multiquadric and the divergence-free trigonometric least-squares post lower KE error ($5.08\%$ and $4.42\%$) than the PI-CON multi-seed mean ($5.71\%$), but both reconstruct the field markedly worse pointwise. Classical interpolation contracts the velocity toward the inter-sensor mean, suppressing total kinetic energy and collapsing the KE ratio toward the DNS value for the wrong reason; the pointwise $u$- and $v$-$L_2$ rows expose this. PI-CON reduces $u$ rel-$L_2$ by $47.2\%$ relative to the strongest fair baseline (trigonometric LSQ, $25.87\%\to13.65\%$) and by $74.2\%$ relative to IDW. Pointwise error, rather than KE alone, separates reconstruction from energy contraction.

The refreshed interpolation sweep also shows strong hyperparameter sensitivity: under the same sensors, over-complete trigonometric LSQ at $k_{\max}=8$ and $12$ and poorly chosen RBF length scales can produce errors orders of magnitude larger than the carefully selected rows in Table~\ref{tab:fair_baselines}. These unstable variants are excluded from the main comparison table and treated as evidence that classical interpolation baselines require explicit validation rather than default trust.

